\begin{textsample}{POS Dim 1 – mt – Score 51.00 – mt/mt\_em\_30.txt}  \label{ex:f1_pos_076}
6.4 Metodologias para o ensino de Ciências da Natureza e suas tecnologias e as aprendizagens significativas O ensino de Ciências da Natureza, no Ensino Médio, deve transpor a ideia de ensino em \textbf{um} viés disciplinar puro, onde cada componente curricular ocupa \textbf{um} espaço definido no currículo em ação.
Nesse contexto, é importante \textbf{salientar} \textbf{que}, na sociedade \textbf{contemporânea}, temos \textbf{uma} juventude imersa em tecnologias onde as informações são fluidas e complementares, isto é, não mais estáticas e \textbf{sólidas} ( BAUMAN, 2001 ).
Sendo assim, a perspectiva de Bauman corrobora com \textbf{uma} \textbf{abordagem} interdisciplinar para o ensino de Ciências, pois é \textbf{imprescindível} \textbf{que} os estudantes percebam \textbf{que} os conhecimentos das “ Ciências da Natureza e suas Tecnologias ” foram produzidos por meio da ação humana e passam por transformações constantes na sociedade \textbf{contemporânea}.
Nesta \textbf{vertente}, a \textbf{divisão} entre as Ciências da Natureza e Ciências Humanas, por exemplo, não é coerente, pois a Física, a Química e a Biologia são também Ciências Humanas, \textbf{visto} \textbf{que} são conhecimentos estabelecidos pelos humanos ( CHASSOT, 2013 ).
Portanto, \textbf{uma} \textbf{abordagem} interdisciplinar, no espaço da sala de aula, torna -se \textbf{um} ponto indissociável da concepção de formação integral dos estudantes, \textbf{uma} vez \textbf{que} os conhecimentos científicos estão \textbf{imbricados} ao processo de desenvolvimento da \textbf{humanidade}.
Nessa conjuntura, é importante apresentar as Ciências aos estudantes como parte de \textbf{um} contexto histórico, levando em conta \textbf{que} nem sempre os \textbf{erros} na ciência levam a resultados desastrosos, \textbf{visto} \textbf{que}, por exemplo, a borracha galvanizada, o teflon e a penicilina foram resultados de \textbf{erros}.
Além disso, citando caso análogo, Camillo Golgi descobriu a coloração com o ósmio, \textbf{uma} técnica para captar detalhes de neurônios visíveis, depois de espirrar o elemento em \textbf{um} tecido cerebral acidentalmente ( KEAN, 2011 ).
Assim, humanizar os conhecimentos científicos, durante as aulas, é \textbf{um} meio para promover a formação integral dos estudantes, de modo \textbf{que} estes sejam protagonistas do seu projeto de vida, capazes de exercer sua cidadania.
Nessa perspectiva, o ensino de Ciências da Natureza, no contexto da BNCC, não pode ignorar \textbf{que} a sociedade do conhecimento é \textbf{baseada} no desenvolvimento de competências cognitivas.
Isso \textbf{posto}, vale \textbf{frisar} \textbf{que}, para Moran : > A escola padronizada, \textbf{que} ensina e avalia a todos de forma igual e \textbf{exige} resultados previsíveis, ignora \textbf{que} a sociedade do conhecimento é \textbf{baseada} em competências cognitivas, pessoais e sociais, \textbf{que} não se adquirem da forma convencional e \textbf{que} \textbf{exigem} proatividade, colaboração, personalização e visão empreendedora. ( 2015, p.16 ).
Ainda de acordo com Moran ( 2010 ), existe \textbf{um} consenso de \textbf{que} estudantes e professores de \textbf{que} as aulas tradicionais ( expositivas ) pautadas na relação teoria-exercício-teoria, em \textbf{uma} perspectiva de \textbf{que} o conhecimento científico se dá por meio de \textbf{memorização} de fórmulas, fatos e \textbf{teorias}, estão ultrapassadas e precisam de mudanças para \textbf{que} os estudantes possam se sentir \textbf{instigados} ao desenvolvimento de competências e habilidades.
Nesse sentido, Barbosa e Moura ( 2013 ) afirmam \textbf{que} muitos professores utilizam metodologias \textbf{que} podem ser consideradas como tipos de metodologias ativas.
Os professores conhecem, mesmo \textbf{que} de forma inconsciente, meios para propiciar ao estudante \textbf{um} processo de aprendizagem mais ativo.
Neste viés, enfatiza -se \textbf{que} : > As metodologias precisam acompanhar os objetivos pretendidos.
Se queremos \textbf{que} os alunos sejam proativos, precisamos adotar metodologias em \textbf{que} os alunos se envolvam em atividades cada vez mais complexas, em \textbf{que} tenham \textbf{que} tomar decisões e avaliar os resultados, com apoio de materiais relevantes.
Se queremos \textbf{que} sejam criativos, eles precisam experimentar \textbf{inúmeras} novas possibilidades de mostrar sua iniciativa. ( MORAN, 2015, p.17 ).
Dessa forma, cresce a necessidade de se trabalhar com as metodologias em \textbf{que} o professor utilize formas de mediação com o propósito de \textbf{que} os estudantes sejam ( e \textbf{que} \textbf{desperte} o desejo de ser ) protagonistas do processo de aprendizagem.
No contexto da BNCC e do DRC/MT–EM, o processo de ensino e aprendizagem planejado com metodologias ativas é de grande relevância, \textbf{visto} \textbf{que} as competências específicas da área são articuladas para \textbf{que} o estudante, a partir delas, se aproprie das competências gerais \textbf{preconizadas} na BNCC.
Assim, é importante a inclusão de metodologias \textbf{que} possibilitem \textbf{que} os estudantes se tornem mais ativos no processo de ensino e aprendizagem, capazes de compreender e interpretar o mundo ( natural, social e tecnológico ) e de transformá -lo com base nos aportes \textbf{teóricos} e processuais das ciências ( BNCC, 2017 ), ou seja, o desenvolvimento do Letramento Científico nos estudantes torna -se \textbf{imprescindível} para a implementação do DRC/MT–EM.
Nessa perspectiva, considera -se \textbf{que} a literatura tem \textbf{uma} vasta referência de metodologias, \textbf{métodos} ou \textbf{abordagens} didáticas \textbf{que} vão ao encontro do \textbf{exposto}.
É importante \textbf{frisar} \textbf{que}, nesse texto, almeja -se apresentar algumas perspectivas de metodologias, consideradas como ativas, \textbf{que} podem auxiliar os professores na mediação das competências e habilidades das Ciências da Natureza e suas Tecnologias.
Dessa forma, para o desenvolvimento das competências das Ciências da Natureza e suas Tecnologias, é de \textbf{suma} importância \textbf{que} os estudantes \textbf{consigam} desenvolver \textbf{uma} visão integrada da interação existente entre a natureza, o conhecimento científico e o desenvolvimento da \textbf{humanidade} e, portanto, \textbf{consigam} mobilizar os conhecimentos construídos no dia a dia, nas aulas de Física, Química e Biologia, para além da observação dos fenômenos naturais, mas na esfera da intervenção social e aplicação no contexto do mundo do trabalho.
Mediante o apresentado no parágrafo anterior, a nível de exemplo, pode -se afirmar \textbf{que} o estudante, ao ler \textbf{uma} reportagem ou ouvir \textbf{uma} notícia nos veículos de comunicação de \textbf{que} o gás utilizado nos fogões é denominado de gás GLP ( Gás Liquefeito do Petróleo ), precisa ter consciência de \textbf{que} o gás GLP é composto de \textbf{uma} mistura de hidrocarbonetos de baixo peso molecular.
Além disso, precisa, sempre \textbf{que} necessário, ser capaz de interpretar estas informações em diferentes contextos, por exemplo, na perspectiva dos processos \textbf{que} ocorreram na natureza ( os processos realizados pelas bactérias, transporte de energia e calor nos sistemas ambientais ) para a formação do petróleo, e também as tecnologias envolvidas em sua extração para produção de combustíveis.
Nessa \textbf{vertente}, para formar \textbf{sujeitos} capazes de utilizar os conhecimentos das Ciências da Natureza e suas Tecnologias de modo crítico e integrado, o Ensino por Investigação apresenta -se como \textbf{uma} possibilidade para o professor no processo de mediação das habilidades da competência específica das Ciências da Natureza e suas Tecnologias.
Assim, por meio do Ensino por Investigação, o estudante será \textbf{instigado} a \textbf{argumentar}, apresentar hipóteses, propor a resolução do problema em questão, analisar os dados e comunicar os resultados.
O professor é peça-chave para favorecer o ensino de maneira \textbf{que} o estudante se torne mais ativo.
Porém, o trabalho do dia a dia com metodologias e estratégias diferenciadas de ensino não é \textbf{tarefa} fácil, mas são essenciais.
Por isso, é \textbf{preciso} pensar em \textbf{um} ensino \textbf{que} proporcione ao estudante \textbf{um} maior protagonismo, seja com materiais de baixo custo, ou feito por meio de \textbf{uma} demonstração investigativa, \textbf{que} ocorre, por exemplo, quando o professor não tem acesso a \textbf{um} lugar específico, como laboratório, e a atividade experimental é realizada em sala de aula.
Ao professor cabe o questionamento, a inclusão de \textbf{problematizações} ou perguntas \textbf{que} façam com \textbf{que} os estudantes possam pensar em \textbf{uma} maneira de como resolver o problema proposto.
A maneira de conduzir a aula é diferenciada, porque o professor faz as perguntas, mas não dá as respostas de \textbf{imediato}.
Isso faz com \textbf{que} os estudantes, além de saírem de suas áreas de conforto, também sejam desafiados a resolverem o problema.
Portanto, o Ensino por Investigação trabalha com a \textbf{vertente} de \textbf{que} o estudante se envolva e se comprometa, tornando -se mais ativo em todo o processo.
Além do Ensino por Investigação, o Estudo de Caso \textbf{configura} -se como \textbf{um} meio para auxiliar os estudantes a fundamentar e defender decisões éticas e responsáveis de modo colaborativo e com o viés interdisciplinar, corroborando para o processo de construção das competências específicas 1, 2 e 3.
Sá e Queiroz ( 2010 ) enfatizam \textbf{que} o \textbf{método} de Estudo de Caso é \textbf{uma} variante do Aprendizado \textbf{Baseado} em Problemas ( ABP ).
Neste sentido, destaca -se \textbf{que} o \textbf{método} de Estudo de Caso possibilita ao professor trabalhar com problemas reais, com o propósito de estimular o desenvolvimento do pensamento crítico, a habilidade de resolução de problemas e a aprendizagem de conceitos ( SÁ E QUEIROZ, 2010 ).
Sendo assim, o \textbf{método} impulsiona os estudantes a buscarem respostas e tomarem decisão sobre a melhor forma de resolver problemas pertinentes para a sociedade.
Com a perspectiva de \textbf{que} o estudante se comporte diferentemente do \textbf{que} está acostumado, tanto o Ensino por Investigação quanto o Estudo de Caso trazem o problema a ser resolvido como ponto-chave para \textbf{que} o estudante inicie \textbf{um} processo de investigação e passe a se familiarizar com normas e valores da comunidade científica.
Assim se evitará \textbf{que} se construa visão deformada sobre ciência, assim como o trabalho do \textbf{cientista}.
Portanto, estas metodologias apresentam -se como possibilidades para serem utilizadas no processo de mediação dos professores para \textbf{desencadear} processos \textbf{que} viabilizem a apropriação das competências específicas das Ciências da Natureza e também as competências gerais da BNCC.
Outro processo \textbf{metodológico} \textbf{que} pode ser \textbf{amplamente} utilizado pelos professores para \textbf{instigar} o processo de investigação científica nos estudantes é a Aprendizagem \textbf{Baseada} em Problemas ( ABP ), \textbf{uma} estratégia de mediação de conhecimentos \textbf{que} se organiza ao redor da investigação de problemas do mundo real ( LOPES \textbf{et} al., 2019 ).
Estudantes e professores se envolvem em analisar, compreender e propor soluções para situações cuidadosamente desenhadas de modo a garantir ao estudante a apropriação das competências específicas 1, 2 e 3, previstas no DRC/MT–EM.
Desta forma, a ABP possibilita \textbf{que} o estudante aplique seus conhecimentos na elaboração de argumentos e \textbf{posicionamentos} frente aos diferentes cenários de \textbf{problematização}, de modo colaborativo e com ética e responsabilidade socioambiental.
Nesse processo de elaboração de argumentos no âmbito da ABP, enfatiza -se a possibilidade de, por meio da mediação do professor, o estudante desenvolver a visão da ciência como \textbf{um} processo contínuo e progressivo \textbf{inerente} às necessidades humanas.
Nessa perspectiva, a ABP é \textbf{um} importante ponto de partida para o planejamento dos professores da área, principalmente para a mediação da competência específica 3, visto \textbf{que} proporciona aos professores meios para motivar os estudantes na mobilização dos conhecimentos científicos na resolução de situações-problema, no contexto de \textbf{uma} visão mais interdisciplinar e transdisciplinar.
Além disso, em síntese, pode -se resumir \textbf{que} a ABP é \textbf{uma} metodologia \textbf{que} se insere no contexto da BNCC como \textbf{uma} importante estratégia de mediação para desenvolver o protagonismo juvenil e \textbf{instigar} o uso de tecnologias digitais de informação e comunicação como meios para o processo de investigação científica.
No Quadro 5, apresenta -se \textbf{uma} síntese das etapas \textbf{que} os professores precisam pensar durante a mediação por meio da ABP.
Quadro 5 - Os “ Sete Saltos ” ( SCHMIDT, 1983 ; FINCO-MAIDAME ; MESQUITA, 2017. ) 1.
Esclarecer frases e conceitos confusos na formulação do problema. 2.
Definir o problema : descrever exatamente \textbf{que} fenômenos devem ser entendidos e explicados. 3.
Chuva de ideias ( Brainstorming ) : usar conhecimentos prévios e senso comum.
Tentar formular o máximo possível de explicações. 4. \textbf{Detalhar} as explicações propostas : tentar construir \textbf{uma} “ \textbf{teoria} ” pessoal, coerente e \textbf{detalhada} dos processos subjacentes aos fenômenos. 5.
Propor temas para a aprendizagem autodirigida. 6.
Preencher as lacunas do conhecimento por meio do estudo individual. 7.
Compartilhar as conclusões com o grupo e integrar os conhecimentos adquiridos em \textbf{uma} explicação adequada dos fenômenos.
Comprovar o \textbf{que} sabe.
Avaliar o processo de construção de conhecimentos.
Dessa forma, é de primordial necessidade o olhar atento do professor, \textbf{que} atuará como mediador ou orientador do processo, considerando o senso comum e os conhecimentos prévios para o desenvolvimento dos novos conhecimentos.
Nesse sentido, a construção do conhecimento na escola \textbf{depende} intrinsecamente da relação de mediação pedagógica estabelecida pelo professor, pois, se o objetivo é formar \textbf{um} cidadão com o desenvolvimento das dez competências gerais, apresentadas na BNCC, é de fundamental importância \textbf{que} a linguagem científica, a criatividade e o protagonismo sejam elementos estruturantes do desenvolvimento do estudante no espaço escolar.
Neste contexto, é relevante reforçar \textbf{que} a ABP é \textbf{uma} metodologia \textbf{que} apresenta variantes, sendo \textbf{que}, para além do Estudo de Caso, conforme já enfatizado em parágrafos anteriores, a aprendizagem \textbf{baseada} em projetos também é \textbf{uma} metodologia variante da ABP.
De forma geral, a aprendizagem \textbf{baseada} em projetos, segundo Bender ( 2014 ), pode ser definida pela utilização de projetos autênticos e realistas, \textbf{baseados} em \textbf{uma} questão, \textbf{tarefa} ou problema \textbf{que} \textbf{instigue} \textbf{um} processo de aprendizagem significativa.
Em síntese, os principais objetivos da aprendizagem \textbf{baseada} em projetos são estimular a motivação dos estudantes, promover \textbf{uma} aprendizagem centrada, incentivar o trabalho em grupo, desenvolver o \textbf{espírito} de iniciativa e criatividade, \textbf{aprimorar} capacidades de comunicação, potencializar o pensamento crítico e proporcionar estratégias para \textbf{que} as habilidades das Ciências da Natureza e suas Tecnologias sejam trabalhadas de forma interdisciplinar ( ALVES ; MOREIRA ; SOUZA, 2007 ). \textbf{Posto} isto, no contexto de proporcionar o desenvolvimento das dez competências gerais apresentadas na BNCC, é de \textbf{suma} relevância articular a aprendizagem \textbf{baseada} em projetos e a \textbf{abordagem} STEAM ( Ciência, Tecnologia, Engenharias, Artes e Matemática, na sigla em inglês – Figura 4 ).
De forma geral, observa -se, por meio da Figura 4, \textbf{que} a \textbf{abordagem} STEAM pode proporcionar aos professores de ciências \textbf{um} planejamento pedagógico mais integrado com as demandas da sociedade \textbf{contemporânea}.
Segundo Bacich e Holanda ( 2020 ), se essa \textbf{abordagem} for associada à metodologia de aprendizagem \textbf{baseada} em projetos, os estudantes poderão ampliar o senso de relevância dos conhecimentos científicos desenvolvidos na educação básica.
Em linhas gerais, alguns elementos são comuns para a aprendizagem \textbf{baseada} em projetos e, consequentemente, também os articuladores dos projetos de STEAM, \textbf{que} são : > [... ] \textbf{um} contexto autêntico capaz de engajar os estudantes, de preferência relacionado a \textbf{um} problema real ; \textbf{uma} sequência de etapas organizadas para a exploração do conhecimento científico e para a produção dos estudantes ; \textbf{um} produto final, geralmente \textbf{um} artefato \textbf{que} permita a aplicação das ideias das engenharias ; e, por fim, a comunicação do projeto, para compartilhar com a comunidade e \textbf{sistematizar} suas aprendizagens. ( BACICH e HOLANDA, 2020, p.19 ) Pensando sob essa forma de organizar o planejamento dos professores de ciências, a \textbf{abordagem} em STEAM, com a metodologia aprendizagem \textbf{baseada} em projetos, precisa ser compreendida como \textbf{um} meio de estimular a criatividade, a investigação científica e a cooperação entre os pares.
Assim, por exemplo, ao propor \textbf{que} os estudantes avaliem o impacto ambiental do descarte inadequado de pneus nas adjacências da escola, e posteriormente solicitar \textbf{que} desenvolvam \textbf{uma} estratégia de intervenção social e ambiental, o fato de os professores da área das Ciências da Natureza debaterem o processo de fabricação dos pneus e também as possíveis doenças tropicais associadas ao descarte inadequado destes, com os estudantes, não se \textbf{configura} como \textbf{uma} \textbf{abordagem} em STEAM ; e os professores de Arte direcionarem os estudantes para a produção de \textbf{um} jardim na escola, a partir da reciclagem dos pneus, pintando -os, como forma de intervenção socioambiental, também não. \textbf{Convém} destacar \textbf{que} o meio de desenvolvimento de \textbf{um} projeto, conforme supracitado, não é inadequado.
Porém, se o objetivo for trabalhar com os estudantes por meio de \textbf{uma} \textbf{abordagem} em STEAM, as ações elencadas \textbf{necessitam} estabelecer \textbf{uma} relação \textbf{dialógica} com os princípios articuladores da \textbf{abordagem} STEAM, \textbf{visto} \textbf{que}, para promovê -la, é \textbf{imprescindível} \textbf{que} os estudantes sejam \textbf{instigados} a analisarem problemas complexos \textbf{que} não possuem \textbf{uma} única solução, mas \textbf{que} possuam contexto autêntico, com viés tecnológico para além do mundo digital, ou seja, como forma de conectar conhecimentos científicos para o desenvolvimento de produtos ( BACICH e HOLANDA, 2020 ).
Além disso, vale \textbf{frisar} \textbf{que}, para formar estudantes aptos à tomada de decisões no campo profissional e pessoal, é importante \textbf{uma} postura de mediação por parte dos professores em \textbf{um} projeto STEAM, ou seja, não direcionar como será o processo de intervenção dos estudantes, mas \textbf{despertar} esses processos criativos e atitudes.
Portanto, a respeito do exemplo mencionado anteriormente, há \textbf{que} se motivar os estudantes a \textbf{um} processo reflexivo em \textbf{que} a compreensão dos conhecimentos científicos sobre as propriedades dos materiais, a análise dos processos de reciclagem e a necessidade de garantir a saúde coletiva sejam \textbf{uma} constante.
Estes são pontos de atenção no planejamento pedagógico dos professores, \textbf{uma} vez \textbf{que} é de \textbf{suma} importância \textbf{que} sejam evitadas, durante a realização do projeto, a \textbf{simples} aplicação de conceitos científicos para a resolução de \textbf{um} problema complexo ou aplicação de roteiros para produção de artefatos.
Nesta perspectiva, enfatiza -se também a necessidade de garantir o alinhamento das competências e habilidades apresentadas neste documento como \textbf{uma} intencionalidade pedagógica do projeto STEAM.
Por último, \textbf{convém} \textbf{salientar} \textbf{que} as proposições dos \textbf{métodos} de ensino, supramencionados, são possibilidades \textbf{que} devem ser expandidas, no cotidiano das escolas, a partir de \textbf{uma} prática pedagógica reflexiva e crítica dos professores das Ciências da Natureza e suas Tecnologias.

% matched lemmas: abordagem, amplamente, aprimorar, argumentar, basear, cientista, configurar, conseguir, contemporâneo, convir, depender, desencadear, despertar, detalhar, dialógico, divisão, erro, espírito, et, exigir, exposto, frisar, humanidade, imbricar, imediato, imprescindível, inerente, instigar, inúmero, memorização, metodológico, método, necessitar, posicionamento, postar, preciso, preconizar, problematização, que, salientar, simples, sistematizar, sujeito, sumo, sólido, tarefa, teoria, teórico, um, ver, vertente
\end{textsample}
